\chapter{Benchmark design and comparative evidence}
A benchmark is a measurement instrument, not a verdict on an architecture.
Before asking whether DOGMA outperforms HermonDNA, specify the behavior measured,
the legal information available, the training opportunity and the population to
which the result is meant to generalize. This chapter turns that specification
into executable scoring, uncertainty calculations, resource controls and
mechanism interventions.

Chapter 22 checked DOGMA's full carry and read-before-write semantics;
Chapter 23 checked masked multi-head attention and gradients.
Chapter 24 established span, reverse-complement, causal-availability and positional
contracts. Chapter 25 audited dependency-closed splits, paired cluster estimates
and architecture-native accounting. Chapter 26 derived a restricted affine scan,
not a general replacement for full DOGMA. Chapters 27--28 added optimization,
resource and artifact checks. These checks enable a benchmark; they do not rank
trained genomic models. The fixtures below are synthetic supplied outcomes,
not hidden training runs or industrial measurements.

\section{Freeze a claim that can fail}
Write the claim before selecting a winner:
\[
 C=(\mathcal T,\mathcal D,M,B,V,\delta_{\min},H).
\]
Here $\mathcal T$ is a task family; $\mathcal D$ fixes records, split/dependence
units and legal access; $M$ is the metric and estimand; $B$ fixes resource
coordinates and tuning opportunity; $V$ names variation sources;
$\delta_{\min}$ is a practically meaningful improvement; and $H$ is the
hypothesis and disconfirming criterion.
\Plate{01-evidence}{A comparative claim depends on task validity, legal access,
resource opportunity and statistical design. A high terminal score cannot repair
an earlier broken link.}{evod29:evidence}

A concrete protocol might ask: under a frozen donor-group split, a declared
training-token ceiling and prespecified tuning budget, does a new mechanism
improve equal-donor mean accuracy by at least $\delta_{\min}$?
Changing to row-weighted accuracy changes the question if donors contribute
unequal numbers of windows. Choosing the threshold after inspecting scores changes
the test. Record amendments and distinguish exploratory from confirmatory results.

\section{Same question, native execution}
DOGMA is non-Transformer; HermonDNA is Transformer-based. Evolutor covers broader
orchestration and research. Their comparison should share data, target, legal access
and scoring interfaces without pretending that native carry and attention/cache
are the same internal object.
\Plate{02-envelope}{The shared envelope fixes evidence semantics; native
state, reset and execution contracts remain distinct. Scores and costs are
joined only after identity and legality checks.}{evod29:envelope}

An offline suffix-aware or reverse-complement representation may be legal for
whole-record classification and illegal for causal next-token prediction.
Neither a good score nor a compact implementation changes that boundary.
For causal tasks, test target shifts, padding and the availability of every feature
at each position. For genomic generalization, audit exact duplicates, reverse
complements, overlaps and shared biological sources before the split.

Begin with a literal oracle and a baseline ladder: majority/random,
GC or length heuristic, simple learned model, established strong baseline,
then proposed mechanism. A dataset solved by GC content has not established
sophisticated sequence reasoning. Failed or null runs belong in the evidence
ledger with configuration and artifact identity; do not invent earlier experiments
or rewrite unsuccessful results as different tasks.

\section{Make the estimand explicit}
For fixed models on the same binary-label test cases, define
\[
 c_i^A=\mathbf1[\hat y_i^A=y_i],\quad
 c_i^B=\mathbf1[\hat y_i^B=y_i],\quad d_i=c_i^A-c_i^B.
\]
Then $\hat\Delta=n^{-1}\sum_i d_i$ is the paired accuracy difference.
\Listing{paired_deltas}
The strict binary schema rejects booleans and ambiguous label encodings.
Multiclass, ranking, calibration and sequence likelihood require their own
metric contracts; accuracy is not a universal proxy for biological usefulness.
For masked likelihood, aggregate loss numerator and valid-token denominator,
not an unweighted mean of differently sized batch means.

The eight-case fixture gives:
\begin{center}\input{\Generated/paired-table.tex}\end{center}
A is correct on seven cases and B on five, hence $\hat\Delta=.25$.
There are three discordant A wins, one discordant B win and four ties.
Pairing retains the common case difficulty. Treating the two accuracy estimates
as independent discards that information.

\section{An exact bootstrap oracle}
For independent identically sampled cases and \emph{fixed} trained models, the
paired percentile bootstrap draws $n$ differences with replacement.
This \Term{paired percentile interval} preserves fixed-model case pairing.
The empirical probabilities here are
\[
 P(D=1)=3/8,\quad P(D=0)=4/8,\quad P(D=-1)=1/8.
\]
Let $H_j(s)$ be the bootstrap sum distribution after $j$ draws:
\[
 H_0(0)=1,\quad
 H_{j+1}(s)=\tfrac38 H_j(s-1)+\tfrac48 H_j(s)+\tfrac18 H_j(s+1).
\]
Exact rational convolution avoids Monte Carlo uncertainty for this tiny problem.
\par\medskip\noindent\begin{minipage}{\linewidth}
\Listing{exact_bootstrap}
\end{minipage}\par\medskip
With the inverse-CDF quantile convention, the nominal 95\% percentile interval is
$[-.25,.75]$. The deterministic 4000-draw bootstrap with seed 29 gives the same
endpoints. Independent Cartesian enumeration checks small fixtures, and an
independent multinomial expansion checks this eight-case distribution.

The calculation is exact for the \emph{empirical resampling law}; its nominal
interval need not have exact frequentist coverage. A one-case sample $d=[1]$
gives a degenerate bootstrap interval $[1,1]$, which plainly does not certify
population certainty. Small samples, boundary effects and dependence demand
caution.

For a different inferential question, an exact paired sign/McNemar-style calculation
conditions on four discordant pairs and assumes exchangeable A/B outcomes under
the null:
\[
 p=\min\left(1,\;2\sum_{j=0}^{\min(3,1)}
                {4\choose j}2^{-4}\right)=.625.
\]
This is not a proof that the methods are equivalent, nor an interchangeable
interpretation of a bootstrap interval. Report assumptions, effect size,
uncertainty and practical threshold rather than collapsing them to a star.

\section{Seeds, groups and crossed dependence}
Training outcomes vary with initialization, data order, augmentation, sampling,
kernel nondeterminism and hyperparameter choices. Bouthillier et al. analyze
these benchmark variance sources \cite{evod29:variance}. A shared integer seed
does not guarantee that two different architectures consume comparable random
streams. Pair shared splits and examples; document architecture-specific randomness.

Five supplied scores $.71,.74,.69,.73,.72$ have mean $.718$ and sample standard
deviation $.0192354$. Report all five, not only the best $.74$.
Their spread is training variation under the declared experiment, not the
uncertainty of a new biological population.

Now define $D_{sg}$ as an A-minus-B group metric for seed pair $s$ and independent
group $g$. Our teaching estimand gives equal weight to both axes:
\[
 \hat\Delta=\frac1{SG}\sum_{s=1}^S\sum_{g=1}^G D_{sg}.
\]
It is not a row-weighted metric and not automatically Chapter 25's ratio of
aggregated numerators/denominators. Unequal group sizes must not silently change it.
\begin{center}\input{\Generated/group-table.tex}\end{center}
The matrix is a hand-specified fixture, not scores from three neural training runs.
Its mean is $1/6$. Resample seed rows and group columns independently, sharing
the same group draw across every chosen row:
\Listing{crossed_bootstrap}
This \Term{crossed bootstrap} retains shared group difficulty across seeds.
It estimates uncertainty for the stipulated exchangeable seed/group design.
It does not include hyperparameter search uncertainty, new-species shift,
arbitrary interactions or a universal coverage guarantee.

For 4000 draws, seed 29, its interval is $[-.5,.75]$.
Resampling seed means only gives $[-.5,.5]$; resampling group means only gives
$[-1/3,2/3]$. These are different variation models, not three interchangeable
confidence levels.
\Plate{03-seeds}{Rows represent training variation; columns represent independent
biological groups. Shared column resampling preserves the crossed design.
The figure shows actual fixture entries, not decorative random dots.}{evod29:seeds}

\subsection{A pseudoreplication counterexample}
If all three seed rows are $[1,1,-1,-1]$, seed means are all zero and a seed-only
interval is $[0,0]$. Yet group uncertainty remains: the crossed interval is
$[-1,1]$ for this fixture. Conversely, identical group effects within each row
can conceal training uncertainty from a group-only analysis.

Duplicating each of two independent effects $+1,-1$ one hundred times creates
200 rows but still only two independent groups. A false IID standard error falls
below $.071$, while the two-unit standard error is $1$. Added windows are not
added independently assigned organisms. The \Term{experimental unit} remains
the independently assigned or sampled group under the declared design. Declare donor, preparation, lineage,
chromosome or family boundaries according to the actual generalization claim.

\section{The resource budget is a vector}
Use a declared set of coordinates, for example
\[
 B=(P,N_{\rm tok},F,M_{\rm peak},M_{\rm state},W_{\rm infer}).
\]
Parameter count, training tokens, forward steps, peak bytes, persistent-state bytes
and inference work are not convertible without an explicit model.
\Listing{within_budget}
Our strict budget type requires nonnegative integer entries and all coordinates.
Unknown peak memory is \emph{unknown}, not zero.
A candidate $(10,10,10,10,10,11)$ fails a ceiling
$(100,100,100,100,100,10)$ despite passing the other five coordinates.
\Plate{04-budget}{Each controlled coordinate has a separate gate.
A parameter advantage cannot compensate for an undeclared inference-work
violation. Hypothetical units are labeled; these values are not measurements.}
{evod29:budget}

Parameter-matched, compute-matched and latency-matched comparisons answer different
questions. Equal forward-step counts can hide different sequence lengths or work per
step. DOGMA carry and HermonDNA KV/attention costs must be accounted natively.
Measured timing requires hardware, warmup, precision, batch/context lengths,
synchronization, percentile reporting and failed-run treatment.
Chapter 25's static arithmetic is not such a timing benchmark.

Our capacity matcher chooses the nearest parameter count with declared tie-breaking;
it does not make every other coordinate fair. A \Term{Pareto comparison} also
requires score noninferiority and no worse controlled costs, with at least one
strict improvement. Otherwise report a tradeoff rather than declaring a universal
winner.

\section{Tuning and task selection change conclusions}
Giving one architecture 200 validation trials and another one default learning rate
changes its opportunity. Freeze search spaces, trial counts or compute budget,
selection metric and stopping rule. Retain failed configurations and consumed work.
Dodge et al. emphasize validation/development reporting beyond test headlines
\cite{evod29:reporting}.

Selection among many noisy results makes the selected maximum optimistic. For
independent estimates with common CDF $F$, $P(\max_{j\leq m}X_j\leq x)=F(x)^m$.
The selected maximum is not an unbiased estimate of a random configuration's
performance. A reserved final test, evaluated after selection freezes, addresses a
different role than the validation set; repeated test-set tuning invalidates that role.

Task selection can even reverse a ranking:
\begin{center}\input{\Generated/lottery-table.tex}\end{center}
A has $.9,.5$ across two tasks and B $.8,.8$. Task one alone favors A;
an equal-weight suite favors B. Dehghani et al. discuss this benchmark-choice
fragility \cite{evod29:lottery}. Freeze task inclusion and aggregation weights.
For multiple confirmatory endpoints, prespecify multiplicity treatment and
distinguish exploratory discoveries. Do not report only tasks on which a
proposed architecture wins.

\section{Mechanism claims require interventions}
An ablation compares a full source model with its mechanism disabled, while checking
changed parameter count, optimization and execution cost. A transplant asks whether
the mechanism helps an interface-compatible recipient. A matched sham checks whether
extra capacity or work explains the gain.
\Plate{05-causal}{Target removal, recipient transplant and a work-matched sham are
distinct interventions. A GC-matched positive/negative pair exposes whether
the detector responds to the intended motif.}{evod29:causal}

Our recipient is a deterministic GC-majority heuristic; the target is a literal
ACG detector. The sham detects ACT using the same full window and character loops.
Both detectors have zero learned parameters and the same static comparison count
for a given length, not a claimed equal hardware latency.
\Listing{intervention}
On the eight-string fixture:
\begin{center}\input{\Generated/mechanism-table.tex}\end{center}
The target improves accuracy from $.375$ to $.75$; the sham does not improve it.
This is a Boolean-program intervention, not a trained DOGMA ablation.
The target is supplied exact label logic, so the example demonstrates the
experiment's causal structure, not learning ability.

The strings ACGTT and ACTGT both have GC count two. Target outputs are $[1,0]$;
sham outputs $[0,1]$. This control defeats a simple GC explanation.
But the OR-combined target still falsely accepts GCGCG: adding a correct positive
detector does not remove the heuristic's false positives.
Exhaustive length-five enumeration checks both detectors against independent
substring membership.

A random-label sanity check must distinguish expected population chance from every
finite empirical score; sufficiently expressive models can also memorize training
labels. The executable negative control here simply inverts the matched pair's
labels: target accuracy falls from one to zero. It is explicitly \emph{not} a
multi-seed shuffled-label training experiment.

\section{Artifacts and claims need separate evidence}
Chapter 28's byte-pinned artifact oracle supports reconstruction for its declared
linear fixture. It is not a full DOGMA/HermonDNA deployment loader.
For a real comparative result, record immutable model/config/tokenizer identities,
dataset and split hashes, inference/reset options, all seeds, selection opportunity,
software/hardware, scoring code and resource observations.

A benchmark card should permit recomputation of every numerator and denominator,
describe group weighting and resampling, state exclusions and failed trials,
and quote the exact bounded claim. A hash authenticates neither biological labels
nor dataset independence; those need their own provenance and audits.
\Plate{06-claims}{Task score, systems performance, mechanism effect and broad
capability are separate claim categories. Extending a conclusion requires new
tasks and controls, not stronger adjectives.}{evod29:claims}

Tokens/s is a systems quantity, not genomic competence. One motif benchmark
is not general genomic reasoning. A mechanism intervention in a hand-coded program
is not proof of a novel learned architecture. Keep negative evidence, deployment
limitations and uncertainty visible when moving from an experimental score to a
professional claim.

\section{Twenty worked exercises}
\begin{enumerate}
\item What is frozen before ranking? \textbf{Answer:} task, access, data/dependence, metric, budgets, variation and claim criterion.
\item Why keep native carry/cache distinct? \textbf{Answer:} their legal state and work semantics differ.
\item Does whole-record classification permit suffix features? \textbf{Answer:} possibly; a causal next-token task does not.
\item Paired fixture accuracy difference? \textbf{Answer:} $7/8-5/8=.25$.
\item Discordant counts? \textbf{Answer:} three A-only wins and one B-only win.
\item Derive $H_{j+1}$. \textbf{Answer:} convolve $H_j$ with empirical masses $3/8,4/8,1/8$ at $1,0,-1$.
\item What is exact about the bootstrap oracle? \textbf{Answer:} its empirical resampling distribution, not nominal coverage.
\item Interpret $[-.25,.75]$. \textbf{Answer:} substantial uncertainty under fixed-model IID-case resampling; not equivalence or proof of A superiority.
\item Conditional sign-test $p$? \textbf{Answer:} $2(1+4)/16=.625$.
\item Five-score mean and sample SD? \textbf{Answer:} $.718$ and $.0192354$.
\item Why not pick only $.74$? \textbf{Answer:} it selects the best trial and hides variation.
\item Crossed matrix estimand? \textbf{Answer:} equal-seed equal-group mean, here $1/6$.
\item Why share column draws across rows? \textbf{Answer:} models share the same biological-group difficulty.
\item Do 100 duplicate windows create 100 independent units? \textbf{Answer:} no.
\item May unmeasured peak bytes be zero? \textbf{Answer:} no; the budget is incomplete.
\item Why does the illustrative candidate fail its budget? \textbf{Answer:} inference work 11 exceeds ceiling 10.
\item Explain the lottery reversal. \textbf{Answer:} A wins task one but loses the frozen two-task mean.
\item What does the ACT sham control? \textbf{Answer:} equal static detector work without the target motif operation, not hardware latency.
\item Why is target accuracy not perfect? \textbf{Answer:} OR preserves GC heuristic false positives.
\item Design a publication gate. \textbf{Rubric:} independent oracle, leakage/access audits,
matched opportunity, named estimand and units, variation/selection plan, unused final
evaluation, intervention/sham, artifact identity, reproducible results and bounded conclusion.
\end{enumerate}

\section{Bridge to runtime contracts}
Chapter 30 defines request identity, state ownership, streaming, cancellation,
deadlines and terminal cleanup. Statistical rigor cannot rescue a runtime that
changes model identity halfway through a request. Correct lifecycle and native
step contracts precede any service-performance comparison.
